
\documentclass[letter, 11pt]{article}

    %%%%%%%%%%%%%%%%%%%%%%%%%%%%%%
    \oddsidemargin  0.0in
    \evensidemargin 0.0in
    \textwidth      6.8in
    \headheight     0.0in
    \topmargin      0.0in
    \textheight=9.0in
    \parindent=0in
    %%%%%%%%%%%%%%%%%%%%%%%%%%%%%%


\renewcommand{\familydefault}{\sfdefault} % Police 

\usepackage{amsthm,amsfonts,amssymb,amsmath, xcolor, enumitem}
\usepackage{graphicx}   %To include graphics
\usepackage[small,compact]{titlesec}

\usepackage[utf8]{inputenc}

\usepackage[french]{babel}

\usepackage{fancyhdr}
\pagestyle{fancy}

\newcommand{\var}{\mathbb{V}\text{ar}} %Symboles programmés pour simplifier le texte



\usepackage{version} % Pour faire un texte troué pour les étudiants
\includeversion{etudiant}\excludeversion{professeur} %version etudiant
%\includeversion{professeur}\excludeversion{etudiants} %version prof
\newcommand{\etudiant}[1]{\begin{etudiants}#1 \end{etudiants}}
\newcommand{\prof}[1]{\begin{professeur}\color{blue} {\bf Solution :} #1 \end{professeur}}


\fancyhead[L]{STT-2100 \\ Régression}
\fancyhead[R]{A.-S. Charest, E. Reny-Nolin, A.-A. Vallée \\ Université Laval}


\begin{document}

\vspace*{0.01 cm}
\begin{center}
{\Large Exercices - Chapitre 2}
\end{center}



\begin{enumerate}[leftmargin=*, wide, label=\bf Exercice \arabic*:]
%QUESTION 1
\item
Soient les données suivantes qui correspondent à la taille ($x_i$ en cm) et au poids ($Y_i$ en kg) de 6 filles de 18 ans (tirées de Morgenthaler, 1997, Exercice 3, p.122) :

\begin{center}
\begin{tabular}{|c|c|c|c|c|c|c|}
\hline
$x_i$ & 170 & 167 & 171 & 181 & 158 & 166 \\\hline
$Y_i$ & 71 & 58 & 56 & 85 & 45 & 53\\\hline
\end{tabular}
\end{center}

\begin{enumerate}[label={\bf \alph*)}] 
\item Déterminez la droite de régression pour les données ci-dessus.
\item  Interprétez la valeur des estimations des paramètres en termes du poids et de la taille des filles de 18 ans de cette population.
\item  Donnez une estimation de la variance des termes d’erreur.
\end{enumerate}

\prof{\begin{enumerate}[label={\bf \alph*)}]
\item  On trouve $n = 6$, $\sum_{i=1}^n x_i = 1013$, $\sum_{i=1}^n x_i^2 = 171\;311$, $\sum_{i=1}^n Y_i = 368$ et $\sum_{i=1}^n x_iY_i = 62\;625$.\\
 On obtient donc que
    $$
    S_{xx} = 171\;311-\frac{1013^2}{6} = 282.83 \text{\quad et \quad} S_{xy} = 62\;625 -\frac{(1013)(368)}{6} = 494.33.
    $$
    De plus $\overline{x}= 168.83$ et $\overline{Y}= 61.33$. \\
		Cela nous conduit aux estimations suivantes $$\hat{\beta}_1= \frac{S_{xy}}{S_{xx}} = \frac{494.33}{282.83} = 1.75$$ $$\hat{\beta}_0 = \overline{Y} - \hat{\beta}_1\overline{x} = 61.33 - 1.75\times168.83 = -234.12.$$
		Le modèle de régression linéaire simple ajusté devient alors $\hat{Y}_i = -234.12 + 1.75x_i$.\\

\item    L'ordonnée à l'origine indique que dans la population des filles de $18$ ans mesurant $0$ cm, le poids moyen est  de $-234.12$ kg. Ce n'est pas une valeur très intéressante à interpréter ici. \\

 La pente indique quant à elle qu'une hausse de la taille de $1$ cm chez une fille de $18$ ans est associée à une augmentation du poids moyen de 1.75 kg. \\

\item L'estimation de la variance des termes d'erreur est
    $$
     \hat\sigma^2=\frac{\sum_{i=1}^n \hat{\varepsilon}_i^2}{n-2} =  \frac{\sum_{i=1}^n(Y_i-\hat{Y}_i)^2}{n-2}
    $$
     qui vaut ici 41.33.
     
    \end{enumerate}}

\item On veut donc ajuster le modèle suivant, dans lequel on force le modèle de régression linéaire à passer par l'origine:
    
    $$Y_i = \beta\,x_i + \varepsilon_i\qquad\text{pour } i=1,\dots,n.$$
    
    Trouver l'estimateur des moindres carrés du paramètre $\beta$ associé à ce modèle.
    
\prof{On veut trouver la valeur de $\beta$ qui minimise la fonction $$L=\sum\limits_{i=1}^n (Y_i-\beta\,x_i)^2.$$ Il faut d'abord annuler sa dérivée par rapport à $\beta$.


$$\begin{array}{rcl}
\dfrac{\partial L}{\partial\beta}&=& 2\,\sum\limits_{i=1}^n (Y_i-\beta\,x_i)(-x_i)\\[6mm]
\dfrac{\partial L}{\partial\beta}=0&\Leftrightarrow &\sum\limits_{i=1}^n (x_iY_i)= \hat\beta \sum\limits_{i=1}^n  x_i^2\\[6mm]
&\Leftrightarrow &\hat\beta = \dfrac{\sum\limits_{i=1}^n (x_iY_i)}{\sum\limits_{i=1}^n  x_i^2}
\end{array}$$

Pour vérifier qu'il s'agit bien d'un minimum, vérifions que la dérivée seconde de $L$ au point $\hat\beta$ est bien positive.

$$\begin{array}{rl}
\dfrac{\partial^2 L}{\partial\beta^2}&= \dfrac{\partial }{\partial\beta}\left[-2\sum\limits_{i=1}^n (x_iY_i) + 2\beta \sum\limits_{i=1}^n x_i^2\right]\\[6mm]
&= 2\sum\limits_{i=1}^n  x_i^2
\end{array}$$

Cette quantité est toujours positive.}

\end{enumerate}

\end{document}

